% Generated draft: requires mathematical and source review.
\documentclass[11pt]{article}
\usepackage[margin=25mm]{geometry}
\usepackage{fontspec}
\setmainfont{Latin Modern Roman}
\usepackage{amsmath,amssymb,mathtools,unicode-math}
\setmathfont{Latin Modern Math}
\usepackage{enumitem,needspace,xurl,hyperref}
\setlength{\parindent}{0pt}
\setlength{\parskip}{5pt}
\setlength{\emergencystretch}{3em}
\newcommand{\N}{\mathbb{N}}
\newcommand{\Z}{\mathbb{Z}}
\newcommand{\Q}{\mathbb{Q}}
\newcommand{\R}{\mathbb{R}}
\newcommand{\C}{\mathbb{C}}
\newcommand{\F}{\mathbb{F}}
\newcommand{\A}{\mathbb{A}}
\newcommand{\PP}{\mathbb{P}}
\newcommand{\E}{\mathbb{E}}
\newcommand{\eps}{\varepsilon}
\newcommand{\dd}{\,\mathrm{d}}
\begin{document}
\section*{UnsolvedMath ID: 700959}
\textbf{Source problem code:} catalogue entry
\subsection{Definitions and mathematical statement}
Let \(\Lambda\) be a totally ordered abelian group, written additively, whose order is invariant under translation. A \(\Lambda\)-metric on a nonempty set \(X\) is a function \(d:X\times X\to\Lambda_{\geq 0}\) satisfying symmetry, the triangle inequality, and \(d(x,y)=0\) if and only if \(x=y\). A segment from \(x\) to \(y\) is the image of an isometric embedding of \([0,d(x,y)]_\Lambda\) with endpoints \(x,y\). A \(\Lambda\)-tree is a \(\Lambda\)-metric space in which any two points are joined by a unique segment, the intersection of two segments with a common endpoint is a segment, and the union of two segments whose intersection consists of their common endpoint is a segment.

An isometric action of a group \(G\) on \(X\) is free on points if
\[
\forall g\in G\ \forall x\in X,\qquad gx=x\Longrightarrow g=1.
\]
Some conventions for free actions on \(\Lambda\)-trees additionally exclude inversions, meaning that no group element exchanges the endpoints of a nondegenerate segment. The supplied record does not specify this convention; it can matter when \(\Lambda\) is not divisible by \(2\).

Use the first-order language of groups \(L_{\mathrm{grp}}=\{\cdot,{}^{-1},1\}\), without predicates for the tree or its action. For any specified class \(\mathcal C\) of groups, its elementary theory is
\[
\operatorname{Th}(\mathcal C)=\{\varphi:\varphi\text{ is an }L_{\mathrm{grp}}\text{-sentence and }G\models\varphi\text{ for every }G\in\mathcal C\}.
\]
This theory is decidable if an algorithm, on every input \(L_{\mathrm{grp}}\)-sentence \(\varphi\), halts and correctly decides whether \(\varphi\in\operatorname{Th}(\mathcal C)\).

\textbf{Problem.} Is \(\operatorname{Th}(\mathcal C)\) decidable, where \(\mathcal C\) is the class of all groups admitting a free isometric action on a \(\Lambda\)-tree? No finite generation or other restriction on the groups is stated.

\textbf{Unresolved scope.} The source does not say whether \(\Lambda\) is fixed or ranges over totally ordered abelian groups. For fixed \(\Lambda\), the class would be
\[
\mathcal C_\Lambda=\{G: G\text{ admits a free isometric action on some }\Lambda\text{-tree}\}.
\]
For varying \(\Lambda\), the class would instead be
\[
\mathcal C_{\mathrm{all}}=\{G:\text{there exist a totally ordered abelian group }\Lambda\text{ and a free isometric action of }G\text{ on a }\Lambda\text{-tree}\}.
\]
The corresponding questions concern \(\operatorname{Th}(\mathcal C_\Lambda)\) and \(\operatorname{Th}(\mathcal C_{\mathrm{all}})\), respectively. The supplied wording does not resolve which scope, or which convention concerning inversions, is intended.
\subsection{Short English statement}
Does an algorithm decide which first-order group sentences hold in every group admitting a free action on a \(\Lambda\)-tree? The source leaves the scope of \(\Lambda\) and the convention concerning inversions unspecified. The registry labels this problem open; that label is supplied source data, not an independently verified status assessment.
\subsection{Sources}
\begin{itemize}
\item Local UnsolvedMath registry, permanent ID \(700959\), exact source code \texttt{catalogue entry}, title \emph{Kourovka Notebook Problem 13.41}. Supplied upstream revision: \texttt{d371ee917ac9647fbb06cd6fc529deec6f3dda82}.
\item \emph{The Kourovka Notebook (Unsolved Problems in Group Theory)}, original problem 13.41, 13th issue (1995); cited in the supplied record through the 21st issue edition (2026), \texttt{arXiv:1401.0300v48}. The supplied collection description names E. I. Khukhro and V. D. Mazurov as editors.
\item Proposed by A. G. Myasnikov and V. N. Remeslennikov, 1995. The registry supplies the status \texttt{open} and an update date of 2026-10-16. The cited edition and mathematical status have not been independently checked for this draft.
\end{itemize}
\end{document}
